\documentclass[a4paper,landscape,8pt]{extarticle}
\usepackage[landscape,margin=0.7cm,top=0.5cm,bottom=0.7cm]{geometry}
\usepackage[utf8]{inputenc}
\usepackage[T1]{fontenc}
\usepackage[ngerman]{babel}
\usepackage{helvet}
\renewcommand{\familydefault}{\sfdefault}
\usepackage{xcolor}
\usepackage{tikz}
\usetikzlibrary{positioning}
\usepackage{enumitem}
\usepackage{multicol}
\usepackage{array}
\usepackage{fancyhdr}
\usepackage{amssymb}

% Farben
\definecolor{searchblue}{HTML}{4285F4}
\definecolor{warnred}{HTML}{C62828}
\definecolor{successgreen}{HTML}{2E7D32}
\definecolor{lightgray}{HTML}{F5F5F5}
\definecolor{bordergray}{HTML}{BBBBBB}

\pagestyle{fancy}
\fancyhf{}
\renewcommand{\headrulewidth}{0pt}
\fancyfoot[C]{\footnotesize Seite \thepage{} von 3}

\setlength{\parindent}{0pt}
\setlength{\parskip}{2pt}

% Kompakte Listen
\setlist{nosep, leftmargin=1em, itemsep=0pt, topsep=1pt}

\begin{document}

% ============================================================================
% SEITE 1: INFORMATIONSBLATT
% ============================================================================

\noindent{\Large\bfseries So finden Sie, was Sie suchen} \hfill
{\footnotesize\textbf{Lernziele:} Gute Suchbegriffe | Werbung erkennen | Ergebnisse bewerten}

\vspace{1mm}
\hrule
\vspace{1mm}

% Drei-Spalten-Layout
\begin{minipage}[t]{0.32\linewidth}
\begin{center}
\colorbox{searchblue!10}{%
\begin{minipage}{0.95\linewidth}
\vspace{2mm}
\begin{center}
{\Large\bfseries\textcolor{searchblue}{1. Gute Suchbegriffe}}\\[1pt]
{\footnotesize\itshape Weniger ist mehr}
\end{center}
\vspace{2mm}
\end{minipage}%
}
\end{center}

\vspace{1mm}

\textbf{Analogie:} Eine Suchmaschine ist wie eine \textbf{Bibliothekarin}. Sie sucht für dich -- aber nur, was du ihr sagst.

\textbf{Schlüsselwörter statt Sätze!}
\begin{itemize}
\item \textcolor{warnred}{\textbf{Schlecht:}} ``Wo finde ich einen guten Hausarzt in Rostock?''
\item \textcolor{successgreen}{\textbf{Gut:}} ``Hausarzt Rostock''
\end{itemize}

\textbf{Spezifisch sein:}
\begin{itemize}
\item \textcolor{warnred}{\textbf{Schlecht:}} ``Kuchen''
\item \textcolor{successgreen}{\textbf{Gut:}} ``Apfelkuchen Rezept einfach''
\end{itemize}

\textbf{Tipps:}
\begin{itemize}
\item 2--4 Wörter reichen meist
\item Keine ganzen Fragen tippen
\item Wichtigste Wörter zuerst
\end{itemize}

\end{minipage}%
\hfill%
\begin{minipage}[t]{0.32\linewidth}
\begin{center}
\colorbox{warnred!10}{%
\begin{minipage}{0.95\linewidth}
\vspace{2mm}
\begin{center}
{\Large\bfseries\textcolor{warnred}{2. Werbung erkennen}}\\[1pt]
{\footnotesize\itshape Nicht alles ist echt}
\end{center}
\vspace{2mm}
\end{minipage}%
}
\end{center}

\vspace{1mm}

\textbf{So erkennen Sie Werbung:}
\begin{itemize}
\item Steht \textbf{``Anzeige''} oder \textbf{``Ad''} dabei
\item Oft \textbf{ganz oben} auf der Seite
\item Manchmal \textbf{leicht anderer Hintergrund}
\end{itemize}

\textbf{Analogie:} Wie Reklame in der Zeitung -- sieht aus wie ein Artikel, ist aber bezahlt.

\textbf{Werbung ist nicht schlecht, aber:}
\begin{itemize}
\item Sie zeigt, wer \textbf{bezahlt} hat
\item Nicht unbedingt das \textbf{beste} Ergebnis
\item Echte Ergebnisse stehen \textbf{darunter}
\end{itemize}

\vspace{2mm}
\fcolorbox{warnred}{warnred!5}{%
\begin{minipage}{0.92\linewidth}
\footnotesize
\textbf{Tipp:} Scrollen Sie ein Stück nach unten -- dort beginnen die echten Suchergebnisse!
\end{minipage}%
}

\end{minipage}%
\hfill%
\begin{minipage}[t]{0.32\linewidth}
\begin{center}
\colorbox{successgreen!10}{%
\begin{minipage}{0.95\linewidth}
\vspace{2mm}
\begin{center}
{\Large\bfseries\textcolor{successgreen}{3. Ergebnisse bewerten}}\\[1pt]
{\footnotesize\itshape Wem kann ich trauen?}
\end{center}
\vspace{2mm}
\end{minipage}%
}
\end{center}

\vspace{1mm}

\textbf{Vertrauenswürdig:}
\begin{itemize}
\item Offizielle Seiten: \textbf{.de}, \textbf{.gov}, \textbf{.org}
\item Bekannte Institutionen:\\Verbraucherzentrale, Stiftung Warentest
\item Zeitungen: tagesschau.de, zeit.de
\end{itemize}

\textbf{Vorsicht bei:}
\begin{itemize}
\item Unbekannten Webseiten
\item Foren und Blogs (nicht geprüft)
\item Seiten mit viel Werbung
\end{itemize}

\textbf{Schnell-Check:}
\begin{itemize}
\item Wer hat das geschrieben?
\item Wann wurde es geschrieben?
\item Klingt es seriös?
\end{itemize}

\end{minipage}

\vspace{1mm}

% Zusammenfassung
\begin{center}
\fcolorbox{bordergray}{lightgray}{%
\begin{minipage}{0.98\linewidth}
\centering
\vspace{1.5mm}
\textbf{Zusammenfassung:}
\textcolor{searchblue}{\textbf{Schlüsselwörter}} statt Sätze ~$|$~
\textcolor{warnred}{\textbf{Werbung}} steht oben (``Anzeige'') ~$|$~
\textcolor{successgreen}{\textbf{Offizielle Seiten}} sind vertrauenswürdiger
\vspace{1.5mm}
\end{minipage}%
}
\end{center}

\newpage

% ============================================================================
% SEITE 2: ÜBUNGEN
% ============================================================================

\begin{center}
{\LARGE\bfseries Übungen: Suchen wie ein Profi}\\[3pt]
{\normalsize Schau dir Seite 1 an und beantworte die Fragen}
\end{center}

\vspace{3mm}
\hrule
\vspace{4mm}

\begin{multicols}{2}

\textbf{\large Teil A: Gute Suchbegriffe}

Welche Suche ist besser? Kreuze an.

\vspace{2mm}

\renewcommand{\arraystretch}{1.8}
\begin{tabular}{@{}p{6cm}p{6cm}c@{}}
1a) ``Wie ist das Wetter morgen?'' & 1b) ``Wetter Rostock morgen'' & $\square$ a ~$\square$ b \\
2a) ``Fahrplan Bahn'' & 2b) ``Zug Rostock Hamburg'' & $\square$ a ~$\square$ b \\
3a) ``Rezept'' & 3b) ``Kartoffelsalat Rezept'' & $\square$ a ~$\square$ b \\
4a) ``Öffnungszeiten IKEA Rostock'' & 4b) ``Wann hat IKEA offen?'' & $\square$ a ~$\square$ b \\
\end{tabular}
\renewcommand{\arraystretch}{1}

\vspace{6mm}

\textbf{\large Teil B: Werbung oder echt?}

Bei welchen Ergebnissen sollten Sie vorsichtig sein?

\vspace{2mm}

\renewcommand{\arraystretch}{1.6}
\begin{tabular}{@{}p{8cm}cc@{}}
& Werbung & Echt \\
1. Ganz oben, ``Anzeige'' steht dabei & $\square$ & $\square$ \\
2. Auf Seite 1, von wikipedia.org & $\square$ & $\square$ \\
3. Leicht gelblicher Hintergrund, ``Ad'' & $\square$ & $\square$ \\
4. Ergebnis von verbraucherzentrale.de & $\square$ & $\square$ \\
\end{tabular}
\renewcommand{\arraystretch}{1}

\vspace{6mm}

\textbf{\large Teil C: Vertrauenswürdigkeit}

Welchen Seiten können Sie eher trauen? Sortieren Sie von 1 (sehr vertrauenswürdig) bis 4 (weniger vertrauenswürdig).

\vspace{2mm}

\renewcommand{\arraystretch}{1.8}
\begin{tabular}{@{}p{7cm}p{2cm}@{}}
a) Ein Forum mit Nutzerbeiträgen & Rang: \dotfill \\
b) verbraucherzentrale.de & Rang: \dotfill \\
c) Eine Seite mit viel blinkender Werbung & Rang: \dotfill \\
d) bundesregierung.de & Rang: \dotfill \\
\end{tabular}
\renewcommand{\arraystretch}{1}

\columnbreak

\textbf{\large Teil D: Praktische Suchen}

Welche Suchbegriffe würden Sie eingeben? Schreiben Sie Ihre Idee auf.

\vspace{3mm}

\textbf{Situation 1:}\\
Sie suchen einen Zahnarzt in Ihrer Stadt.

\textit{Ihre Suchbegriffe:}
\vspace{6mm}

\rule{\linewidth}{0.4pt}

\vspace{5mm}

\textbf{Situation 2:}\\
Sie möchten wissen, wie spät der Supermarkt heute schließt.

\textit{Ihre Suchbegriffe:}
\vspace{6mm}

\rule{\linewidth}{0.4pt}

\vspace{5mm}

\textbf{Situation 3:}\\
Sie wollen wissen, welche Nebenwirkungen ein Medikament hat.

\textit{Ihre Suchbegriffe:}
\vspace{6mm}

\rule{\linewidth}{0.4pt}

\vspace{5mm}

\textbf{Situation 4:}\\
Sie suchen ein einfaches Rezept für Pfannkuchen.

\textit{Ihre Suchbegriffe:}
\vspace{6mm}

\rule{\linewidth}{0.4pt}

\end{multicols}

\newpage

% ============================================================================
% SEITE 3: CHECKLISTE
% ============================================================================

\begin{center}
{\LARGE\bfseries Checkliste: Das habe ich verstanden}\\[3pt]
{\normalsize Geh die Liste durch und hake ab, was du sicher weißt}
\end{center}

\vspace{3mm}
\hrule
\vspace{6mm}

% Drei Boxen nebeneinander
\begin{minipage}[t]{0.31\linewidth}
\begin{center}
\fcolorbox{searchblue}{searchblue!8}{%
\begin{minipage}{0.92\linewidth}
\vspace{3mm}
\begin{center}
{\large\bfseries\textcolor{searchblue}{Suchbegriffe}}
\end{center}
\vspace{2mm}
\begin{itemize}[label=$\square$, itemsep=5pt]
\item Ich verwende \textbf{Schlüsselwörter} statt ganzer Sätze
\item Ich weiß: \textbf{2--4 Wörter} reichen meist
\item Ich kann \textbf{spezifisch} suchen
\item Ich setze das Wichtigste \textbf{zuerst}
\end{itemize}
\vspace{3mm}
\end{minipage}%
}
\end{center}
\end{minipage}%
\hfill%
\begin{minipage}[t]{0.31\linewidth}
\begin{center}
\fcolorbox{warnred}{warnred!8}{%
\begin{minipage}{0.92\linewidth}
\vspace{3mm}
\begin{center}
{\large\bfseries\textcolor{warnred}{Werbung erkennen}}
\end{center}
\vspace{2mm}
\begin{itemize}[label=$\square$, itemsep=5pt]
\item Ich weiß, dass \textbf{``Anzeige''} bezahlte Werbung bedeutet
\item Ich schaue \textbf{unter} die Werbung
\item Ich weiß, dass Werbung nicht das \textbf{beste} Ergebnis ist
\item Ich scrolle \textbf{nach unten} für echte Ergebnisse
\end{itemize}
\vspace{3mm}
\end{minipage}%
}
\end{center}
\end{minipage}%
\hfill%
\begin{minipage}[t]{0.31\linewidth}
\begin{center}
\fcolorbox{successgreen}{successgreen!8}{%
\begin{minipage}{0.92\linewidth}
\vspace{3mm}
\begin{center}
{\large\bfseries\textcolor{successgreen}{Ergebnisse bewerten}}
\end{center}
\vspace{2mm}
\begin{itemize}[label=$\square$, itemsep=5pt]
\item Ich kenne \textbf{vertrauenswürdige} Seiten
\item Ich bin \textbf{vorsichtig} bei Foren und Blogs
\item Ich frage mich: \textbf{Wer} hat das geschrieben?
\item Ich achte auf das \textbf{Datum}
\end{itemize}
\vspace{3mm}
\end{minipage}%
}
\end{center}
\end{minipage}

\vspace{8mm}

% Gesamtverständnis
\begin{center}
\fcolorbox{black}{lightgray}{%
\begin{minipage}{0.95\linewidth}
\vspace{4mm}
\begin{center}
{\large\bfseries Gesamtverständnis}
\end{center}
\vspace{3mm}
\begin{multicols}{2}
\begin{itemize}[label=$\square$, itemsep=6pt]
\item Ich kann \textbf{effektiv suchen} mit wenigen Wörtern
\item Ich erkenne \textbf{Werbung} in den Suchergebnissen
\item Ich weiß, welchen Seiten ich \textbf{trauen} kann
\item Ich kann \textbf{praktische Suchen} durchführen (Arzt, Fahrplan, Rezept)
\item Ich fühle mich \textbf{sicherer} bei der Internetsuche
\end{itemize}
\end{multicols}
\vspace{3mm}
\end{minipage}%
}
\end{center}

\vspace{8mm}

% Notfall-Notizen
\begin{center}
\fcolorbox{bordergray}{white}{%
\begin{minipage}{0.95\linewidth}
\vspace{4mm}
\begin{center}
{\large\bfseries Meine Such-Notizen}\\[3pt]
{\footnotesize (Fülle das hier aus und bewahre es auf!)}
\end{center}
\vspace{4mm}
\renewcommand{\arraystretch}{2}
\begin{tabular}{@{}p{0.45\linewidth}p{0.45\linewidth}@{}}
\textbf{Meine meistgesuchten Themen:} & \textbf{Vertrauenswürdige Seiten, die ich kenne:} \\
\dotfill & \dotfill \\[3mm]
\dotfill & \dotfill \\[3mm]
\dotfill & \dotfill \\
\end{tabular}
\renewcommand{\arraystretch}{1}
\vspace{4mm}
\end{minipage}%
}
\end{center}

\vfill

\begin{center}
\footnotesize\textit{Stand: \today{} -- Bei Zweifeln an einem Suchergebnis: Lieber auf bekannte Seiten gehen!}
\end{center}

\end{document}
